\documentclass[11pt]{article}

\usepackage[T1]{fontenc}
\usepackage{amsmath,amssymb,amsthm}
\usepackage[margin=1in]{geometry}
\usepackage[colorlinks=true,linkcolor=blue,urlcolor=blue]{hyperref}

\newtheorem{theorem}{Theorem}
\newtheorem{lemma}[theorem]{Lemma}
\newtheorem{proposition}[theorem]{Proposition}
\newtheorem{conjecture}[theorem]{Conjecture}
\theoremstyle{remark}
\newtheorem{remark}[theorem]{Remark}

\newcommand{\R}{\mathbb R}
\newcommand{\Pp}{\mathbb P}
\newcommand{\norm}[1]{\left\|#1\right\|}
\newcommand{\ip}[2]{\left\langle #1,#2\right\rangle}
\newcommand{\weakto}{\rightharpoonup}
\newcommand{\weakstarto}{\overset{*}{\rightharpoonup}}
\newcommand{\diver}{\mathop{\rm div}}
\newcommand{\supp}{\mathop{\rm supp}}
\setlength{\emergencystretch}{2em}

\title{The Full q4 Defect as a Nonlinear Entrance\\
from Frequency Infinity}
\author{John Lawson and Codex}
\date{25 July 2026}

\begin{document}
\maketitle

\begin{abstract}
This round sharpens the conditional energy-efficient exact q4 branch.
The preceding full-defect construction splits the velocity as \(u=a+b\),
where \(a\) is a projected backward Oseen entrance carrying the entire
downward kinetic-energy jump and \(b\) converges strongly to the terminal
trace.  On reversing terminal time and changing both signs,
\[
 v(\tau):=-a(T-\tau),\qquad c(\tau):=-b(T-\tau),
\]
the missing component satisfies the forward nonlinear equation
\[
 \partial_\tau v-\nu\Delta v+
 \Pp(((v+c)\cdot\nabla)v)=0.
\]
It has weak vector trace zero, positive limiting energy \(d\), and initial
energy measure equal to the positive terminal defect \(\vartheta\).
Moreover, it retains fixed energy on every selected q4 amplitude layer and
\(\norm{v(\tau_j)}_{L^{3,\infty}}\gtrsim m_j\).

Exact local and spectral relative-energy identities are derived.  They do
not close the branch: the two positive defect measures cancel in the
relative density, and the surviving commutator has no sign or
clock-integrable majorant supplied by the current hypotheses.  An exact
Zeno ledger realises order-one nonlinear clock capacity with finite
unweighted but infinite amplitude-weighted action.  Thus the next gate is
a genuinely new zero-trace nonlinear-entrance, boundary-energy,
inverse-cascade, or quantitative compactness theorem.  Every result is
conditional on the repository's q4 chain and awaits external review.  No
Clay alternative is proved.
\end{abstract}

\section{Epistemic status and inherited hypotheses}

Let \(u\) be a smooth finite-energy solution of
\begin{equation}
\label{eq:nse}
 \partial_tu-\nu\Delta u+(u\cdot\nabla)u+\nabla p=0,
 \qquad \diver u=0,
\end{equation}
on \(\R^3\times(0,T)\), with a Leray--Hopf continuation through the
putative first singular time \(T\).  This report works only inside the
conditional energy-efficient exact q4 cell established in preceding
rounds.  In particular, write
\begin{equation}
\label{eq:defect}
 E_-:=\lim_{t\uparrow T}\norm{u(t)}_2^2,\qquad
 u_*:=u(T),\qquad
 d:=E_--\norm{u_*}_2^2>0.
\end{equation}

The full-defect projected-Oseen theorem supplies a solenoidal field
\[
 a\in L^\infty(0,T;L^2_\sigma)
 \cap L^2(0,T;\dot H^1_\sigma)
\]
satisfying, in pressure form,
\begin{equation}
\label{eq:adjoint}
 \partial_ta+\nu\Delta a+(u\cdot\nabla)a+\nabla q=0,
 \qquad \diver a=0,
\end{equation}
and
\begin{align}
\label{eq:pair}
 \ip{u(t)}{a(t)}&=d,\\
\label{eq:aweak}
 a(t)&\weakto0\quad\hbox{in }L^2\quad(t\uparrow T),\\
\label{eq:aenergy}
 \norm{a(t)}_2^2+
 2\nu\int_t^T\norm{\nabla a(s)}_2^2\,ds&=d.
\end{align}
If
\begin{equation}
\label{eq:bdef}
 b:=u-a,
\end{equation}
then
\begin{equation}
\label{eq:bstrong}
 b(t)\longrightarrow u_*
 \quad\hbox{strongly in }L^2
 \quad(t\uparrow T).
\end{equation}

There is an exact first-record sequence \(t_j\uparrow T\), with
\(\tau_j:=T-t_j\), amplitudes \(m_j\to\infty\), and layers
\begin{equation}
\label{eq:layers}
 A_j:=\{\alpha_j<|u(t_j)|\leq2\alpha_j\}
\end{equation}
such that, for fixed \(e_0>0\),
\begin{equation}
\label{eq:layer-data}
 \int_{A_j}|u(t_j)|^2\,dx\geq e_0,\qquad
 R_j:=|A_j|^{1/3}\asymp m_j^{-2}.
\end{equation}
Along this sequence,
\begin{equation}
\label{eq:measure-input}
 |a(t_j)|^2\,dx\weakstarto\vartheta,
\end{equation}
where
\begin{equation}
\label{eq:theta}
 \vartheta\geq0,\qquad
 \vartheta(\R^3)=d,\qquad
 \supp\vartheta\subset\sigma,\qquad
 \mathcal H^1(\sigma)=0.
\end{equation}
These are hypotheses of this round, not conclusions about arbitrary
Clay-admissible data.

\section{Terminal reversal produces a nonlinear entrance}

\begin{theorem}[Forward perturbed Navier--Stokes entrance]
\label{thm:reverse}
For \(\tau=T-t\), define
\begin{equation}
\label{eq:vc}
 v(\tau):=-a(T-\tau),\qquad
 c(\tau):=-b(T-\tau).
\end{equation}
Then
\begin{equation}
\label{eq:v-projected}
 \partial_\tau v-\nu\Delta v+
 \Pp(((v+c)\cdot\nabla)v)=0,
 \qquad \diver v=\diver c=0.
\end{equation}
Equivalently,
\begin{equation}
\label{eq:v-pressure}
 \partial_\tau v-\nu\Delta v+
 ((v+c)\cdot\nabla)v+\nabla Q=0,
\end{equation}
where \(Q(\tau):=q(T-\tau)\) and
\begin{equation}
\label{eq:pressure}
 -\Delta Q=
 \partial_i\partial_\ell\bigl((v_\ell+c_\ell)v_i\bigr).
\end{equation}
At the initial boundary,
\begin{align}
\label{eq:vweak}
 v(\tau)&\weakto0,\\
\label{eq:cstrong}
 c(\tau)&\longrightarrow-u_*
 \quad\hbox{strongly in }L^2,\\
\label{eq:venergy}
 \norm{v(\tau)}_2^2+
 2\nu\int_0^\tau\norm{\nabla v(s)}_2^2\,ds&=d.
\end{align}
\end{theorem}

\begin{proof}
Since \(u=a+b\), equation \eqref{eq:adjoint} reads
\[
 \partial_ta+\nu\Delta a+
 ((a+b)\cdot\nabla)a+\nabla q=0.
\]
The definitions \eqref{eq:vc} imply
\[
 u(T-\tau)=-(v+c),\qquad
 \partial_ta(T-\tau)=\partial_\tau v,\qquad
 \Delta a(T-\tau)=-\Delta v.
\]
Substitution gives \eqref{eq:v-pressure}; Leray projection gives
\eqref{eq:v-projected}.  Taking divergence of
\eqref{eq:v-pressure}, using both divergence constraints, gives
\eqref{eq:pressure}.

Equations \eqref{eq:vweak} and \eqref{eq:cstrong} are
\eqref{eq:aweak} and \eqref{eq:bstrong} after the change of variables.
The same change of variables in \eqref{eq:aenergy} proves
\eqref{eq:venergy}.
\end{proof}

\begin{remark}[The precise boundary failure]
For \(0<\varepsilon<\tau\), equation \eqref{eq:venergy} gives
\begin{equation}
\label{eq:positive-energy}
 \norm{v(\tau)}_2^2+
 2\nu\int_\varepsilon^\tau\norm{\nabla v(s)}_2^2\,ds
 =\norm{v(\varepsilon)}_2^2.
\end{equation}
Thus \(v\) obeys the ordinary dissipative equality away from zero, but
\[
 v(\tau)\weakto0,\qquad
 \lim_{\tau\downarrow0}\norm{v(\tau)}_2^2=d>0.
\]
Extending the ordinary energy inequality to \(\tau=0\) with datum zero
would force \(v\equiv0\).  No such extension has been proved under the
present supercritical drift regularity.
\end{remark}

\section{The positive initial energy measure}

Put \(e_v:=|v|^2/2\) and
\begin{equation}
\label{eq:flux}
 F_v:=(v+c)e_v+Qv-\nu\nabla e_v.
\end{equation}

\begin{proposition}[Local energy law with measure trace]
\label{prop:local-energy}
For every positive time,
\begin{equation}
\label{eq:local-energy}
 \partial_\tau e_v+\diver F_v=-\nu|\nabla v|^2.
\end{equation}
Moreover, for every \(\phi\in C_c^\infty(\R^3)\) and fixed \(\tau>0\),
\begin{equation}
\label{eq:measure-law}
 \frac12\int_{\R^3}\phi\,d\vartheta
 =
 \int_{\R^3}\phi e_v(\tau)\,dx
 +\nu\int_0^\tau\!\!\int_{\R^3}
       \phi|\nabla v|^2\,dx\,ds
 -\int_0^\tau\!\!\int_{\R^3}
       F_v\cdot\nabla\phi\,dx\,ds.
\end{equation}
Equivalently, as spatial distributions,
\begin{equation}
\label{eq:measure-distribution}
 \frac12\vartheta
 =
 e_v(\tau)\,dx+
 \nu\int_0^\tau|\nabla v(s)|^2\,ds\,dx+
 \diver\int_0^\tau F_v(s)\,ds.
\end{equation}
\end{proposition}

\begin{proof}
Dot \eqref{eq:v-pressure} with \(v\).  The divergence constraints give
\[
 v\cdot((v+c)\cdot\nabla v)
 =\diver((v+c)e_v),\qquad
 v\cdot\nabla Q=\diver(Qv).
\]
The diffusion identity
\[
 -\nu v\cdot\Delta v
 =-\nu\Delta e_v+\nu|\nabla v|^2
\]
then proves \eqref{eq:local-energy}.

The energy-class interpolation
\[
 L^\infty_\tau L^2_x\cap L^2_\tau L^6_x
 \subset L^{10/3}_{\tau,x}
\]
holds for \(v\), \(c\), and \(v+c\).  Hence the cubic flux belongs
locally to \(L^{10/9}_{\tau,x}\).  Equation \eqref{eq:pressure} and
Calderon--Zygmund boundedness give \(Q\in L^{5/3}_{\tau,x}\), so
\(Qv\in L^{10/9}_{\tau,x}\).  Finally,
\(\nabla e_v=v_i\nabla v_i\) is locally integrable.  Thus all flux
integrals below are absolutely continuous at the lower time boundary.

Integrate \eqref{eq:local-energy} from
\(\varepsilon=\tau_j\) to \(\tau\) against \(\phi\).  By
\eqref{eq:measure-input},
\[
 |v(\tau_j)|^2\,dx=|a(t_j)|^2\,dx
 \weakstarto\vartheta.
\]
Letting \(j\to\infty\) proves \eqref{eq:measure-law}.
The definition of distributional divergence gives
\eqref{eq:measure-distribution}.
\end{proof}

\begin{remark}
The weak vector datum is zero, but the scalar energy datum is
\(\vartheta/2\).  A successful boundary theorem must therefore use more
than the vector barycentre of the escaping oscillations.
\end{remark}

\section{The entrance inherits the q4 amplitude layer}

\begin{lemma}[Finite-volume weak-\(L^3\) estimate]
\label{lem:lorentz}
For every measurable \(E\subset\R^3\) of finite measure and every
\(f\in L^{3,\infty}(\R^3)\),
\begin{equation}
\label{eq:lorentz-set}
 \int_E|f|^2\,dx
 \leq C|E|^{1/3}\norm{f}_{L^{3,\infty}}^2.
\end{equation}
\end{lemma}

\begin{proof}
Let \(K:=\norm{f}_{L^{3,\infty}}\).  The distribution function obeys
\[
 |\{|f|>\lambda\}|\leq K^3\lambda^{-3}.
\]
Using the layer-cake formula on \(E\) and splitting at
\(\lambda_0:=K|E|^{-1/3}\) gives
\[
 \int_E|f|^2
 \leq
 2\int_0^{\lambda_0}\lambda|E|\,d\lambda+
 2\int_{\lambda_0}^{\infty}\lambda K^3\lambda^{-3}\,d\lambda
 \leq C K^2|E|^{1/3}.
\]
\end{proof}

\begin{theorem}[q4 amplitude inheritance]
\label{thm:amplitude}
For all sufficiently large \(j\),
\begin{equation}
\label{eq:carrier-energy}
 \int_{A_j}|v(\tau_j)|^2\,dx\geq\frac{e_0}{4},
\end{equation}
and
\begin{equation}
\label{eq:weak-l3}
 \norm{v(\tau_j)}_{L^{3,\infty}}\gtrsim m_j.
\end{equation}
\end{theorem}

\begin{proof}
Since \(b(t_j)\to u_*\) strongly in \(L^2\),
\[
 \int_{A_j}|b(t_j)|^2\,dx
 \leq
 2\norm{b(t_j)-u_*}_2^2+
 2\int_{A_j}|u_*|^2\,dx.
\]
The first term tends to zero.  The second does too because
\(|A_j|=R_j^3\to0\) and \(|u_*|^2\in L^1\).  Therefore
\[
 \norm{a(t_j)}_{L^2(A_j)}
 \geq
 \norm{u(t_j)}_{L^2(A_j)}
 -\norm{b(t_j)}_{L^2(A_j)}
 \geq\sqrt{e_0}-o(1).
\]
Since \(v(\tau_j)=-a(t_j)\), this proves
\eqref{eq:carrier-energy}.  Lemma~\ref{lem:lorentz},
\eqref{eq:layer-data}, and \(|A_j|^{1/3}=R_j\asymp m_j^{-2}\)
then imply
\[
 \frac{e_0}{4}
 \lesssim R_j\norm{v(\tau_j)}_{L^{3,\infty}}^2,
\]
which is \eqref{eq:weak-l3}.
\end{proof}

\section{Exact relative-current cancellation}

Subtracting \eqref{eq:adjoint} from \eqref{eq:nse} gives
\begin{equation}
\label{eq:b-equation}
 \partial_tb-\nu\Delta b+(u\cdot\nabla)b+\nabla(p-q)
 =2\nu\Delta a.
\end{equation}
Define
\begin{equation}
\label{eq:r}
 r:=a\cdot b
\end{equation}
and, using repeated-index summation,
\begin{equation}
\label{eq:relative-current}
 \begin{aligned}
 \mathcal J_k^{\rm rel}:={}&
 u_kr+(p-q)a_k+qb_k\\
 &+\nu\left(
 u_i\partial_ka_i-a_i\partial_ku_i
 -2a_i\partial_ka_i
 \right).
 \end{aligned}
\end{equation}

\begin{theorem}[Local relative-current law]
\label{thm:relative-local}
For every preterminal time,
\begin{equation}
\label{eq:relative-law}
 \partial_tr+\diver\mathcal J^{\rm rel}
 =-2\nu|\nabla a|^2.
\end{equation}
At the terminal boundary,
\begin{equation}
\label{eq:r-zero}
 r(t)\,dx\weakto0
\end{equation}
as a locally finite signed measure, while
\begin{equation}
\label{eq:r-total}
 \int_{\R^3}r(t)\,dx
 =2\nu\int_t^T\norm{\nabla a(s)}_2^2\,ds.
\end{equation}
Consequently, for every \(\phi\in C_c^\infty(\R^3)\),
\begin{equation}
\label{eq:relative-representation}
 2\nu\int_t^T\!\!\int_{\R^3}\phi|\nabla a|^2\,dx\,ds
 =
 \int_{\R^3}\phi\,a(t)\cdot b(t)\,dx
 +\int_t^T\!\!\int_{\R^3}
 \mathcal J^{\rm rel}\cdot\nabla\phi\,dx\,ds.
\end{equation}
\end{theorem}

\begin{proof}
The primal--adjoint cross identity is
\[
 \partial_t(u\cdot a)+\diver J^{\rm cross}=0,
\]
where
\[
 J_k^{\rm cross}
 =
 u_k(u\cdot a)+pa_k+qu_k
 -\nu(\partial_ku_i\,a_i-u_i\partial_ka_i).
\]
Directly from \eqref{eq:adjoint},
\[
 \partial_t|a|^2+
 \diver\left(u|a|^2+2qa+2\nu a_i\nabla a_i\right)
 =2\nu|\nabla a|^2.
\]
Subtract the second identity from the first and use
\(u\cdot a-|a|^2=a\cdot b\).  The resulting flux is precisely
\eqref{eq:relative-current}, proving \eqref{eq:relative-law}.

For compactly supported \(\phi\),
\[
 \int\phi\,a\cdot b
 =
 \ip{a}{\phi(b-u_*)}+\ip{a}{\phi u_*}.
\]
The first term tends to zero by \eqref{eq:bstrong} and the boundedness
of \(a\) in \(L^2\); the second tends to zero by \eqref{eq:aweak}.
This proves \eqref{eq:r-zero}.  Globally,
\[
 \int r=\ip{a}{u}-\norm{a}_2^2
 =d-\norm{a}_2^2,
\]
so \eqref{eq:aenergy} proves \eqref{eq:r-total}.  Finally, integrate
\eqref{eq:relative-law} from \(t\) to \(T\), test against \(\phi\),
and use \eqref{eq:r-zero} to obtain
\eqref{eq:relative-representation}.
\end{proof}

\begin{remark}[Why positivity disappears]
Along the q4 records, both \(u\cdot a\,dx\) and
\(|a|^2\,dx\) converge to the same positive measure
\(\vartheta\).  Their difference \(r\,dx\) therefore has zero terminal
measure.  The remaining local current in
\eqref{eq:relative-representation} has no known sign.
\end{remark}

\section{Exact pressure-free spectral balance}

Let \(B_N=\psi(D/N)\), where \(0\leq\psi\leq1\) is smooth, radial,
nonincreasing in frequency, equal to one on the unit ball, and supported
in the ball of radius two.  Thus \(B_N\) is self-adjoint and commutes with
derivatives and the Leray projection.  Put
\[
 D_u:=u\cdot\nabla,\qquad
 [D_u,B_N]:=D_uB_N-B_ND_u
\]
and
\begin{equation}
\label{eq:RN}
 R_N(t):=\ip{b(t)}{B_Na(t)}.
\end{equation}

\begin{theorem}[Spectral relative-energy identity]
\label{thm:spectral}
For every \(N>0\),
\begin{equation}
\label{eq:RN-derivative}
 \frac{d}{dt}R_N
 =
 \ip{[D_u,B_N]b}{a}
 -2\nu\ip{\nabla a}{B_N\nabla a}.
\end{equation}
Moreover,
\begin{equation}
\label{eq:RN-zero}
 \sup_{N>0}|R_N(t)|\longrightarrow0
 \qquad(t\uparrow T),
\end{equation}
and hence
\begin{equation}
\label{eq:RN-integrated}
 \begin{aligned}
 R_N(t)
 ={}&
 2\nu\int_t^T\ip{\nabla a}{B_N\nabla a}\,ds\\
 &-\int_t^T\ip{[D_u,B_N]b}{a}\,ds.
 \end{aligned}
\end{equation}
\end{theorem}

\begin{proof}
The exact low-pass cross identity is
\begin{equation}
\label{eq:cross-low}
 \frac{d}{dt}\ip{u}{B_Na}
 =\ip{[D_u,B_N]u}{a}.
\end{equation}
Indeed, the two viscous terms and both pressures cancel, while
the skew-adjointness of \(D_u\) converts the transport terms into the
commutator.  Similarly,
\begin{equation}
\label{eq:a-low}
 \frac{d}{dt}\ip{a}{B_Na}
 =
 2\nu\ip{\nabla a}{B_N\nabla a}
 +\ip{[D_u,B_N]a}{a}.
\end{equation}
Since
\[
 R_N=\ip{u}{B_Na}-\ip{a}{B_Na},
\]
subtracting \eqref{eq:a-low} from \eqref{eq:cross-low} proves
\eqref{eq:RN-derivative}.

For completeness, \eqref{eq:RN-zero} also follows directly from the
strong--weak terminal decomposition.  Write
\[
 R_N=\ip{b-u_*}{B_Na}+\ip{B_Nu_*}{a}.
\]
The first term tends to zero uniformly in \(N\).  The set
\[
 \{B_Nu_*:0<N<\infty\}\cup\{0,u_*\}
\]
is norm compact in \(L^2\): the multiplier orbit is norm-continuous,
with limits zero as \(N\downarrow0\) and \(u_*\) as \(N\uparrow\infty\).
Weak convergence of the bounded family \(a(t)\) is uniform on this
compact set, so the second term also tends uniformly to zero.
Integrating \eqref{eq:RN-derivative} to \(T\) gives
\eqref{eq:RN-integrated}.  With \(B_N\) replaced by the identity, the
commutator vanishes and the formula reduces to
\eqref{eq:r-total}; there is no additional global positive charge.
\end{proof}

\begin{lemma}[Available commutator estimates]
\label{lem:commutator}
For almost every time set
\[
 M:=\norm{u}_{L^{3,\infty}},\qquad
 Y:=\norm{\nabla a}_2,\qquad
 Z:=\norm{\nabla b}_2.
\]
Then
\begin{equation}
\label{eq:MZY}
 \left|\ip{[D_u,B_N]b}{a}\right|
 \lesssim MZY.
\end{equation}
If
\[
 A:=\sup_{t<T}\norm{a(t)}_2,\qquad
 B:=\sup_{t<T}\norm{b(t)}_2,
\]
then also
\begin{equation}
\label{eq:bernstein-comm}
 \left|\ip{[D_u,B_N]b}{a}\right|
 \lesssim BNMY+ABN^2M.
\end{equation}
\end{lemma}

\begin{proof}
By skew-adjointness, Lorentz Holder, and the endpoint Sobolev embedding,
\[
 \begin{aligned}
 |\ip{D_uB_Nb}{a}|
 &=
 |\ip{B_Nb}{D_ua}|\\
 &\lesssim
 \norm{B_Nb}_{L^{6,2}}\,M\,Y
 \lesssim MZY.
 \end{aligned}
\]
Likewise,
\[
 \begin{aligned}
 |\ip{B_ND_ub}{a}|
 &=
 |\ip{b}{D_uB_Na}|\\
 &\lesssim
 \norm{b}_{L^{6,2}}\,M\,\norm{\nabla B_Na}_2
 \lesssim MZY.
 \end{aligned}
\]
This proves \eqref{eq:MZY}.  Alternatively,
\[
 \norm{B_Nb}_{L^{6,2}}\lesssim NB,\qquad
 \norm{\nabla B_Na}_{L^{6,2}}\lesssim N^2A.
\]
Use the first bound in the first commutator term and pair the second
term in \(L^2\) to obtain \eqref{eq:bernstein-comm}.
\end{proof}

\begin{remark}[Exact limitation]
The energy class gives \(Y^2,Z^2\in L^1_t\), while the q4 chain has
unbounded \(M\) and forces \(\int_t^TMY^2=\infty\).  This does
\emph{not} prove that the relative commutator diverges.  It proves that
\eqref{eq:MZY} and \eqref{eq:bernstein-comm} furnish no
clock-uniform integrable majorant.  Strong convergence of \(b\) makes
its terminal path norm compact and its high-frequency \(L^2\) tail
uniformly small, but supplies no rate after multiplication by \(N\) or
\(N^2\), nor a high-frequency gradient modulus.
\end{remark}

\section{The sharp Zeno survivor}

\begin{proposition}[Exact q4 clock ledger]
\label{prop:zeno}
Let
\begin{equation}
\label{eq:zeno-def}
 m_j:=2^{2j},\qquad
 \delta_j:=\frac{2^{-11j}}{j},\qquad
 \Lambda_j:=2^{14j/3}j^{2/3}.
\end{equation}
Then
\begin{equation}
\label{eq:capacity}
 m_j^2\Lambda_j^{3/2}\delta_j=1.
\end{equation}
If \(q_j:=\delta_j\Lambda_j^2\), then
\begin{equation}
\label{eq:qj}
 q_j=2^{-5j/3}j^{1/3},\qquad
 \sum_jq_j<\infty,
\end{equation}
but
\begin{equation}
\label{eq:weighted-qj}
 m_jq_j=2^{j/3}j^{1/3},\qquad
 \sum_jm_jq_j=\infty.
\end{equation}
\end{proposition}

\begin{proof}
The powers in \eqref{eq:capacity} are
\[
 2^{4j}\left(2^{14j/3}j^{2/3}\right)^{3/2}
 \frac{2^{-11j}}j
 =
 2^{4j+7j-11j}j^{1-1}=1.
\]
Squaring \(\Lambda_j\) gives
\[
 q_j=
 \frac{2^{-11j}}j\,2^{28j/3}j^{4/3}
 =2^{-5j/3}j^{1/3},
\]
which is summable.  Multiplication by \(m_j=2^{2j}\) gives
\eqref{eq:weighted-qj}.
\end{proof}

\begin{proposition}[Hilbert-frequency scalar survivor]
\label{prop:hilbert}
Let \(\{e_j\}\) be an orthonormal sequence in an infinite-dimensional
Hilbert space, assigned frequencies \(\Lambda_j\).  Put
\[
 D_j:=\sum_{k\geq j}q_k,\qquad
 A_j:=(d-2\nu D_j)^{1/2}
\]
for sufficiently large \(j\), and define the record states
\[
 v_j:=A_je_j.
\]
Choose a fixed \(c_*\) orthogonal to every \(e_j\), and set
\[
 c_j:=c_*+\beta_je_j,\qquad
 \beta_j:=\frac{2\nu D_j}{A_j}.
\]
Then
\begin{align}
\label{eq:hilbert-weak}
 v_j&\weakto0,&
 \norm{v_j}_2^2+2\nu D_j&=d,\\
\label{eq:hilbert-strong}
 c_j&\longrightarrow c_*\quad\hbox{strongly},&
 \ip{v_j}{c_j}&=2\nu D_j.
\end{align}
\end{proposition}

\begin{proof}
Since \(q_j\) is summable, \(D_j\downarrow0\), so
\(A_j\to\sqrt d\).  Bounded coefficients multiplying an orthonormal
sequence converge weakly to zero, proving the first part of
\eqref{eq:hilbert-weak}; the second is the definition of \(A_j\).
Also \(\beta_j\to0\), which proves strong convergence in
\eqref{eq:hilbert-strong}, while
\(\ip{v_j}{c_j}=A_j\beta_j=2\nu D_j\).
\end{proof}

\begin{remark}
Proposition~\ref{prop:hilbert} realises, at the scalar record level, a
weak-zero positive-energy boundary, a strongly convergent remainder,
the exact relative overlap, and frequency escape.  Together with
Proposition~\ref{prop:zeno}, it defeats any contradiction using only
those scalar budgets.  It supplies neither spatial fields nor a pressure,
a continuous path, a local current, or a Navier--Stokes solution.  Genuine
one-event smooth triad packets from the preceding sharpness round also
vary the trajectory from event to event.  A same-trajectory PDE theorem
remains essential.
\end{remark}

\section{What changed in this round}

\subsection*{Robust conditional findings, subject to external review}

\begin{enumerate}
 \item The entire missing kinetic component is a forward nonlinear
 perturbed-Navier--Stokes field after terminal reversal, not merely a
 remote dual detector.
 \item Its vector trace is weakly zero, its drift trace is strongly
 prescribed, and its scalar initial energy is the positive measure
 \(\vartheta/2\).
 \item The nonlinear entrance inherits fixed carrier energy and the
 diverging q4 weak-\(L^3\) amplitude.
 \item The local relative-current law
 \eqref{eq:relative-law} and spectral identity
 \eqref{eq:RN-integrated} are exact and pressure-consistent.
 \item Relative subtraction cancels the two copies of
 \(\vartheta\); it does not reveal a second fixed positive reservoir.
 \item Present commutator estimates expose a quantitative compactness
 gap rather than a contradiction.
 \item The Zeno and Hilbert ledgers satisfy all currently used scalar
 energy, clock, overlap, and action constraints without being a PDE
 counterexample.
\end{enumerate}

\subsection*{Closed shortcut}

Strong \(L^2\) convergence of \(b\) alone does not make the relative
commutator summable.  It gives qualitative compactness without the
clock-relative decay rate needed to absorb the factors \(N\) or \(N^2\).
Likewise, subtracting the common positive defect removes its positivity
from the relative balance.

\subsection*{Open obligations}

\begin{enumerate}
 \item Prove a zero-trace nonlinear-entrance theorem under the exact q4
 same-trajectory hypotheses, forcing \(d=0\).
 \item Equivalently, extend \eqref{eq:positive-energy} to an energy
 inequality at \(\tau=0\), using the strong-trace drift and the q4
 coupling rather than assuming strong initial convergence.
 \item Force a non-summable inverse-cascade, residence,
 projective-motion, capacity, or pressure charge.
 \item Obtain a quantitative clock-relative compactness modulus for
 \(c\), or an integrable estimate for the commutator in
 \eqref{eq:RN-integrated}.
 \item Exclude \(\vartheta\) using its \(\mathcal H^1\)-null support or
 the coupled pressure equation.
 \item Subject the inherited full-defect alignment and every theorem in
 this report to independent external mathematical review.
 \item Treat slower clocks, divergent normalised energy, R3B, and the
 other Clay alternatives separately.
\end{enumerate}

\begin{conjecture}[No q4 nonlinear entrance from infinity]
\label{conj:no-entrance}
Suppose
\[
 v,c\in
 L^\infty(0,\tau_0;L^2_\sigma)
 \cap L^2(0,\tau_0;\dot H^1_\sigma)
\]
satisfy \eqref{eq:v-projected}, with
\[
 v(\tau)\weakto0,\qquad
 c(\tau)\to c_0\ \hbox{strongly in }L^2,\qquad
 \lim_{\tau\downarrow0}\norm{v(\tau)}_2^2=d.
\]
If \(v+c\) is the exact reversal of one smooth Navier--Stokes trajectory
in the energy-efficient q4 first-record cell, and \(v\) carries the q4
layers \eqref{eq:carrier-energy}--\eqref{eq:weak-l3}, then \(d=0\).
\end{conjecture}

\begin{remark}[Scope of the conjecture]
The same-trajectory and q4 clauses are essential.  The Hilbert ledger
shows that the scalar boundary data alone do not force \(d=0\), and this
report makes no assertion that generic weak solutions with rough drift
have the proposed boundary rigidity.  Conjecture~\ref{conj:no-entrance}
is not proved.
\end{remark}

\section*{Conclusion}

The conditional q4 problem has been converted into a sharper and more
recognisable PDE question: can a nonlinear viscous field come down from
frequency infinity with zero weak vector trace but a positive singular
energy trace, while its coupled drift converges strongly?  The conversion,
measure law, amplitude inheritance, and relative identities are proved
under the inherited branch assumptions.  The decisive boundary theorem is
not.  Accordingly, the Clay problem remains unsolved.

\end{document}
